\documentclass[10pt,a4paper]{amsart}
\usepackage[margin=19mm]{geometry}
\usepackage{amsmath,amssymb,mathtools}
\usepackage[scr=rsfs,cal=euler]{mathalfa}
\usepackage{xfrac}
\usepackage[unicode,hidelinks,bookmarks=false]{hyperref}
\newcommand{\ie}{i.e.\ }
\newcommand{\cf}{cf.\ }
\newcommand{\resp}{respectively\ }
\newcommand{\Subsection}{\subsection}
\providecommand{\nobreakdash}{\nobreak\hbox{-}}
\setlength{\emergencystretch}{2em}
\multlinegap0pt
\usepackage{bm}
\usepackage{bbm}
\usepackage{dsfont}
\usepackage{xspace}
\usepackage{cancel}
\usepackage{mathtools}
\usepackage{cleveref}
\usepackage{amsthm}
\makeatletter
\@ifundefined{claim}{\newtheorem{claim}{}}{}
\@ifundefined{theorem}{\newtheorem{theorem}{Theorem}}{}
\@ifundefined{lemma}{\newtheorem{lemma}[theorem]{Lemma}}{}
\makeatother
\title[Splitter Theorems for Graph Immersions]{Splitter Theorems for Graph Immersions}
\author{Matt DeVos, Mahdieh Malekian}
\date{Source: arXiv:1808.06569v3 (CC BY 4.0)}
\begin{document}
\maketitle

\noindent\emph{Source and licence.} Splitter Theorems for Graph Immersions. Version arXiv:1808.06569v3 (CC BY 4.0), distributed under the Creative Commons Attribution 4.0 International licence, as recorded in the arXiv licence metadata for this exact version and shown on the abstract page https://arxiv.org/abs/1808.06569v3. Full rights notice: licenses/arxiv-1808.06569-CC-BY-4.0.txt.\par\smallskip

\noindent\textit{Editorial scope.} Counted unit: the Lemma whose source label is \texttt{lem10} in the article (source0.tex lines 716--719, source label \texttt{lem10}), delivered together with the article's own complete original proof of it (source0.tex lines 721--747, one \texttt{proof} environment of 27 lines). No mathematics is written, repaired or re-derived: statement and proof are the article's own text line for line. Required context retained uncounted: sp-thm2 (lines 617--625); lem9 (lines 632--635); 3econdel (lines 701--704). Every definition, display and label of those lines is retained unchanged. Omitted from this package and not redistributed: 1--616, 626--631, 636--700, 705--715, 720--720, 748--860. Nothing inside a retained excerpt is omitted or altered: every retained body is delivered exactly as the article's lines are written, so no omission receipt of any kind accompanies this package. Citations: the retained excerpts carry no citation command at all, so no bibliography entry of the article is transcribed and no reference is redistributed. Reference adaptation: none was needed and none was made; every reference inside the delivered unit resolves inside it, so no unresolved reference or citation is produced. Printed slips kept literal: no printed word, symbol, hypothesis or number of the counted statement or of its proof was altered. Layout and class: the article's own class is \texttt{article} with packages including xcolor, amsmath, bbold, xy, graphicx, enumitem, marginnote, amsthm, hyperref, cleveref; the wrapper is the lane's pinned \texttt{amsart} editorial wrapper at 10pt on A4, which restates the theorem-like environments the retained text needs and adds the article's own macro definitions that the retained text needs (none), with extra wrapper packages (bm, bbm, dsfont, xspace, cancel, mathtools, cleveref). The delivered numbering is the unit's own. Assets: no retained span contains a figure environment, a graphics-inclusion command, a PDF-page inclusion, a table or a TikZ picture; no image file is copied into this package, the package declares no asset dependency and no image file is redistributed with it. Licence: version arXiv:1808.06569v3 of this article is distributed under the Creative Commons Attribution 4.0 International licence, as recorded in the arXiv licence metadata for this exact version and shown on the abstract page https://arxiv.org/abs/1808.06569v3. A full rights notice is delivered at licenses/arxiv-1808.06569-CC-BY-4.0.txt. Characters: every retained excerpt is pure ASCII, so the delivered engine, XeLaTeX with its default Latin Modern text font, reports no missing character.
\begin{theorem}
\label{sp-thm2}
Let $G$ and $H$ be $3$-edge-connected and internally $4$-edge-connected loopless graphs, with $G \succ H$. Assume that $|V(H)|\ge 2$, and $(G, H) \ncong (Q_3, K_4), (Q_3, K_2^3)$. Then there exists an operation taking $G$ to $G'$ such that $G'$ is $3$-edge-connected, internally $4$-edge-connected, and $G' \succeq   H$, where an operation is either
\begin{itemize}
\item deleting an edge,
\item splitting at a vertex of degree at least 4,
\end{itemize}
each followed by iteratively deleting any loops, and suppressing vertices of degree $2$.
\end{theorem}

\begin{lemma}
\label{lem9}
Let $G$ and $H$ be as in \cref{sp-thm2}. Suppose there is an edge $e \in E(G)$ such that $G \setminus e$ has an $H$-immersion. If $(G, H) \ncong (Q_3, K_4), (Q_3, K_2^3)$, then a good operation exists.
\end{lemma}

\begin{lemma}
\label{3econdel}
Let $H$ be a $3$-edge-connected graph, and let $Y$ be a minimal subset of $V(H)$ such that $\delta(Y)$ is a nontrivial $3$-edge-cut in $H$. Then for every edge $e$ in $H[Y]$, $H\setminus e$ is internally $3$-edge-connected.
\end{lemma}

\begin{lemma}
\label{lem10}
Let $G$ and $H$ be as in \cref{sp-thm2}. Assume that there is a split at a vertex $v \in V(G)$ preserving an $H$-immersion. If $(G, H) \ncong (Q_3, K_4), (Q_3, K_2^3)$, then a good operation exists.
\end{lemma}

\begin{proof}
Let $uvw$ be the $2$-edge-path that is to be split. Note if $d(v)=3$, then deleting the edge incident to $v$ other than $vu,vw$ preserves the $H$-immersion. Hence, by \cref{lem9} we are done. Also, observe that if a split is done at a vertex of degree at least four, the resulting graph is $3$-edge-connected. Therefore, we only need to look into the case where splitting off $uvw$ destroys internal $4$-edge-connectivity. So, it must be the case that $uv, vw$ contribute to some 4- or 5-edge-cut $\delta (X)=\{ uv, wv, x_1y_1, x_2y_2 (, x_3y_3): u, w, x_i \in X \}$, where $|X|,|X^c|\ge 2$. We now split the analysis into cases depending on $d(X)$.
\begin{claim}
\label{splitdX4}
If $d(X)=4$, a good operation exists.
\end{claim}
%%%%%%%%%%%%%%%%%%
Since $H$ is $3$-edge-connected and splitting $uvw$ creates a 2-edge-cut, all terminals of $H$ lie on one side of the cut. Also, since $G$ is $3$-edge-connected, each side of the cut contains an edge lying in it, i.e. $E(G[X]), E(G[X^c])\neq \emptyset$.

First, suppose all terminals of $H$ are in $X$. Observe that if we can delete an edge in $G[X^c]$ in a way that  the connectivity between $y_1$ and $y_2$ in $G[X^c]$ is preserved, an $H$-immersion is present in the resulting graph. We propose to delete an edge $e \in E(G[X^c])$ and claim that deleting $e$ preserves the $H$-immersion. It suffices to show $e$ is not a cut-edge in $G[X^c]$ separating $y_1, y_2$. For a contradiction, suppose $e=\delta(Y)$ separates $y_1,y_2$ in $G[X^c]$, where $y_1\in Y$. We may also assume, without loss of generality, that $v \in Y$. Then $\delta_G(Y^c)$ would be a $2$-edge-cut in $G$, a contradiction. Therefore, we can delete $e$ using \cref{lem9}.

Next, suppose all terminals of $H$ are in $X^c$. Similar to the previous case, if we modify $X$ in a way that preserves the connectivity between $x_1, x_2$ in $G[X]$, then an $H$-immersion is sure to exist in the resulting graph. Again, we propose to delete an edge $e \in E(G[X])$ and claim that deleting $e$ preserves the $H$-immersion. It suffices to show $e$ is not a cut-edge in $G[X]$ separating $x_1, x_2$. For a contradiction, suppose $e=\delta(Y)$ separates $x_1,x_2$ in $G[X]$, where $x_1\in Y$. Note $3$-edge-connectivity of $G$ implies that $\delta(Y)$ separates $u, w$ as well. We may assume, without loss of generality, that $u\in Y, w\in Y^c$. Then $d_G(Y)=d_G(Y^c)=3$, thus it follows from internal $4$-edge-connectivity of $G$ that $|Y|=|Y^c|=1$ and $Y=\{u=x_1\}, Y^c=\{w=x_2\}$. Therefore, $X$ consists of two vertices $u,w$ of degree three, and thus deleting $uw$ preserves the $H$-immersion.
%%%%%%%%%%%%%%%%%%

\begin{claim}
If $d(X)=5$, a good operation exists.
\end{claim}
By the internal edge-connectivity of $H$, all terminals of $H$ but possibly one, lie on one side of the cut. First, suppose that most terminals of $H$ are in $X$. Observe that if $X^c$ is modified in a way that preserves the presence of three edge-disjoint paths form a vertex in it to $X$ not using $uv, vw$, the presence of $H$-immersion is guaranteed. Next, suppose that most terminals of $H$ are in $X^c$. In this case, if we manage to modify $X$ in a way that preserves the presence of three edge-disjoint paths form a vertex in it to $X^c$ covering $\delta(X)\setminus \{uv,vw\}$, the presence of $H$-immersion is guaranteed. 

Now, we claim the aforementioned modifications are possible whether most terminals are in $X$ or in $X^c$. In any case, let $G'$ be the graph resulting from splitting off $uvw$, followed by suppressing $v$ in case $d_{G}(v)=4$. We denote the edge created by splitting $uvw$ by $e'$. Note, by \ref{splitdX4}, we may assume $G'$ is 3-edge-connected.

Take an arbitrary nontrivial $3$-edge-cut $\delta_{G'}(Y)$ in $G'$. Observe that $\delta_G(Y)$ must have been a $5$-edge-cut in $G$, which both edges of the split $2$-path $uvw$ contributed to. So, in particular, $e'$ lies either completely in $Y$ or in $Y^c$. Also, there must be an edge other than $e'$ in $G'[Y]$. It is because $3$-edge-connectivity of $G$ implies $6\le \sum_{y \in Y} d_G(y)=d_G(Y)+2 e_G(G[Y])=5+2 e_G(G[Y])$. Thus $e_G(G[Y])>0$, and so there is an edge $\neq e'$ in $G'[Y]$.

Now, let $Z$ denote the side of $\delta(X)$ containing most terminals of $H$ (, so $Z=X$ or $X^c$). We will show that there is an edge lying in $Z^c$ which we could delete, while preserving an $H$- immersion. Since $\delta_{G'}(Z)$ is a nontrivial $3$-edge-cut, we may choose a minimal $Y\subseteq Z^c$ such that $\delta_{G'}(Y)$ is a nontrivial $3$-edge-cut.

It is argued above that there exists an edge $e\neq e'$ in $G'[Y]$. We claim deletion of $e$ preserves the $H$-immersion. It is because it follows from \cref{3econdel} that $G'\setminus e$ is internally $3$-edge-connected. Now, $3$-edge-connectivity of $G'$ and $d_{G'}(Y)=3$ imply that $G'[Y]\setminus e$ has a vertex of degree at least three. Therefore, there exists in $G'\setminus e$ three edge-disjoint paths from such a vertex to $Z$. Observe that since these set of paths cover $\delta(Z)$, deletion of $e$ from $G$ preserves the presence of $H$-immersion. We now can use \cref{lem9} to delete $e$ from $G$.
\end{proof}

\end{document}
